% Lösungen Einheit 3 – Klammern auflösen (zusammenbau v1.2)
\kbabschnitt{Klammern auflösen}

\lbloes{13.}{a)}{ein Plus}
\lbloes{}{b)}{ein Minus}
\lbloes{}{c)}{die Zahl $5$ als Faktor}
\lbloes{}{d)}{$-6$ als Faktor}
\lbloes{14.}{a)}{$5a + 2b - 1$}
\lbloes{}{b)}{$5a + 2b - 3$}
\lbloes{}{c)}{$5x + 15$}
\lbloes{}{d)}{$10$}
\lbloes{}{e)}{$5a - 2b - 4$}
\lbloes{}{f)}{$7a - 2b + 5 - c$}
\lbloes{}{g)}{$-2x - 10$}
\lbloes{}{h)}{$7x + 12$}
\lbloes{}{i)}{$4x + 23$}
\lbloes{15.}{a)}{Nicht stimmen können die Rechnungen 2 und 3. (2): Die Zahl vor der Klammer muss mit jedem Glied malgenommen werden, hinten muss ein Vielfaches von 3 stehen. (3): Das Minus vor der Klammer dreht alle Vorzeichen, aus $-6$ muss $+6$ werden.}
\lbloes{}{b)}{(1) wahr, denn Plus vor der Klammer ändert kein Vorzeichen. (2) falsch, z. B. $10 - (x + 1) = 10 - x - 1$: aus $+x$ wird $-x$. (3) wahr, denn so sagt es das Distributivgesetz.}
\lbloes{}{c)}{$5x + 10$}
\lbloes{}{d)}{$216 + 9x$}
